% UTF-8; compile with XeLaTeX or Tectonic (XeTeX).
\documentclass[UTF8,fontset=fandol,11pt,a4paper]{ctexart}
\usepackage[margin=24mm,headheight=15pt]{geometry}
\usepackage{amsmath,amssymb,mathtools,bm}
\usepackage{graphicx,xcolor,booktabs,array,tabularx,longtable}
\usepackage{caption,placeins,fancyhdr,enumitem}
\usepackage[unicode,colorlinks,linkcolor=blue!50!black,urlcolor=blue!50!black]{hyperref}
\usepackage{bookmark}
\hypersetup{pdftitle={从部分受控交换到QRAM查询：线路可行性与相干脉冲验证},pdfauthor={TCG}}
\graphicspath{{images/}}
\setlength{\parindent}{2em}
\setlength{\parskip}{.25em}
\setlength{\emergencystretch}{2em}
\linespread{1.23}
\captionsetup{font=small,labelfont=bf,labelsep=quad}
\setlist{nosep,topsep=.5em,leftmargin=2.2em}
\renewcommand{\arraystretch}{1.23}
\pagestyle{fancy}\fancyhf{}
\fancyhead[L]{\small QRAM与分布式Grover}
\fancyhead[R]{\small 线路可行性与相干脉冲验证}
\fancyfoot[C]{\thepage}
\newcommand{\ket}[1]{\lvert #1\rangle}
\newcommand{\bra}[1]{\langle #1\rvert}
\newcommand{\proj}[1]{\ket{#1}\!\bra{#1}}
\newcommand{\tr}{\operatorname{Tr}}
\newcommand{\diag}{\operatorname{diag}}
\newcommand{\CZ}{\mathrm{CZ}}
\newcommand{\Fe}{F_{\mathrm e}}
\newcommand{\seq}{\;\longrightarrow\;}
\newcommand{\code}[1]{\texttt{\detokenize{#1}}}
\newcommand{\datafig}[3]{\begin{figure}[htbp]\centering
\includegraphics[width=\linewidth,height=.65\textheight,keepaspectratio]{#1.pdf}
\caption{#2}\label{#3}\end{figure}}
% Exported from saved coherent data, not measured hardware.
\newcommand{\IdealError}{8.483\times10^{-15}}
\newcommand{\RfiveGrover}{98.9007\%}
\newcommand{\RfiveFe}{97.9840\%}
\newcommand{\RfiveGHZ}{98.4727\%}
\newcommand{\RsixGrover}{98.7237\%}
\newcommand{\RsixFe}{98.3080\%}
\newcommand{\RsixGHZ}{98.6344\%}
\newcommand{\RfivePairPop}{99.7340\%}
\newcommand{\RfiveNoPairPop}{84.8445\%}
\newcommand{\RfiveLocalFe}{99.5699\%}
\newcommand{\RfiveNoPairFe}{95.7375\%}

\begin{document}
\hypersetup{pageanchor=false}
\begin{titlepage}\centering
\vspace*{2cm}
{\Huge\bfseries 从部分受控交换到QRAM查询\par}
\vspace{.8cm}{\LARGE 线路可行性与相干脉冲验证\par}
\vspace{.6cm}{\Large R5/R6节点、固定数据库访问与单轮Grover\par}
\vspace{1.5cm}
\begin{minipage}{.88\linewidth}\large
本报告依次回答：查询目标是什么；现有门集能否构造节点；部分节点的入口约束在整树中是否成立；真实脉冲的高能级相位怎样影响组合；最终查询和Grover是否保留正确相干关系。

公式使用原生\LaTeX{}排版。数据图来自已保存的相干仿真，参数、完整单门序列和图像随工程提供。
\end{minipage}
\vfill{\large 2026年11月2日\par}
\vspace{.5cm}{\small 模拟研究，不是真机结果。主结果不含耗散、附加随机门错误或测量误差。\par}
\end{titlepage}
\hypersetup{pageanchor=true}
\section*{先读这一页：结论与阅读顺序}
\addcontentsline{toc}{section}{先读这一页：结论与阅读顺序}
在规定的空端口入口下，R5与R6均能由CZ、01半门、01与12的$\pi$门及虚拟参考帧实现正确路由。按本文顺序连接地址加载、下行、固定存储访问与相干恢复，理想原生矩阵检查通过全部16种数据库；完整查询矩阵最大振幅误差为$\IdealError$。

冻结局域脉冲导出的整机相干基线中，R5/R6四目标平均原始Grover成功率为\RfiveGrover{}与\RsixGrover{}，16数据库平均查询$\Fe$为\RfiveFe{}与\RsixFe{}。这支持参考线路的相干组合可行性，不等于真实17位器件已有同等性能。
\begin{center}\small\begin{tabularx}{\linewidth}{@{}p{.19\linewidth}XX@{}}\toprule
证据 & 回答的问题 & 本报告范围\\\midrule
理想原生门 & 逻辑结构是否正确 & 展开节点与接口门，检查完整复数查询矩阵\\
局域四能级脉冲 & 高能级相位与人口能否控制 & 同一门库的下行、上行、实际往返与消融\\
整树模块过程 & 节点能否组合成查询和算法 & 13个量子站点加4个基态存储条件变量\\
完整真实器件 & 所有共享端口与跨片门是否成立 & 尚待真实拓扑和各节点门库补齐\\\bottomrule
\end{tabularx}\end{center}
阅读顺序是“目标与门集---节点机制---完整查询---脉冲模型---逐层验证---算法闭环---资源与边界”。逐门相位放在附录，避免打断主线。所有时间箭头均从左向右，矩阵乘积仍由右侧先作用。
\tableofcontents\clearpage
\section{先定义任务：正确不只是读出一个比特}\label{sec:task}
两地址位$A_0,A_1$组成$a\in\{00,01,10,11\}$，$D$是外部数据位。四个存储位按地址次序制备为经典基态$\ket m_M=\ket{m_{00}m_{01}m_{10}m_{11}}$。临时路由寄存器记为$\mathcal R$，目标查询为
\begin{equation}\label{eq:querygoal}
 Q_m\ket a_A\ket d_D\ket0_{\mathcal R}\ket m_M
 =\ket a_A\ket{d\oplus m_a}_D\ket0_{\mathcal R}\ket m_M.
\end{equation}
这必须对任意输入的线性叠加成立。地址分支之间不能多出未知相位，临时树也不能留下“走过哪条路”的记录，否则计算基读出正确仍可能无法完成干涉算法。

因此本报告同时检查：人口是否按目标置换；复数相位是否正确；临时寄存器是否恢复。人口真值表只覆盖第一项。本轮关闭$T_1/T_2$耗散、附加随机门错误、准备翻转、读出混淆和有限shots重建，但保留冻结脉冲的非理想旋转、高能级相位与相干泄露。所有密度矩阵由传播直接得到，不是假装已完成实验QST。

\section{从允许门集到部分router}\label{sec:native}
\subsection{单门的相位与原生展开}
四线次序统一为$C,I,L,R$，依次是控制、输入、左输出、右输出。目标跃迁$\ket j,\ket{j+1}$内采用
\begin{equation}\label{eq:rotation}
 R_{j,j+1}(\theta,\phi)=\begin{pmatrix}
 \cos(\theta/2)&-i e^{-i\phi}\sin(\theta/2)\\
 -i e^{i\phi}\sin(\theta/2)&\cos(\theta/2)
 \end{pmatrix}.
\end{equation}
$X_{01}(\phi)=R_{01}(\pi,\phi)$，$X_{01}/2(\phi)=R_{01}(\pi/2,\phi)$，$X_{12}(\phi)=R_{12}(\pi,\phi)$。未写轴时名义相位为0，$\pm Y_{01}/2$对应$\phi=\pm\pi/2$。

$\pi$门在被驱动块中是$-iX$；若第三能级旁路，$-i$未必是整条线路的全局相位，不能粗略把所有X当成无相位置换。

只有CZ是原生两比特门。CNOT时间顺序为
\begin{equation}\label{eq:cnot}
 -Y_{01}^{(b)}/2\seq\CZ(a,b)\seq+Y_{01}^{(b)}/2,
 \qquad C_{a\to b}=R_y^{(b)}(\pi/2)\CZ_{ab}R_y^{(b)}(-\pi/2).
\end{equation}
实际CZ和单门的原生相位补偿由编译器累计到帧，不能只取裸波形而省去补偿。

接收端已知为0时，两CNOT搬运为
\begin{equation}\label{eq:transfer}
 T_{a\to b}: C_{a\to b}\seq C_{b\to a},\qquad
 T_{a\to b}\ket{\psi,0}=\ket{0,\psi},
\end{equation}
用2个CZ和4个半门。若两端未知，完整01空间SWAP应为
\begin{equation}\label{eq:swap}
 S_{ab}: C_{a\to b}\seq C_{b\to a}\seq C_{a\to b},
\end{equation}
用3个CZ和6个半门；它不表示任意qutrit交换。外部H也实际展开：
\begin{equation}\label{eq:h}
 +Y_{01}/2\seq X_{01},\qquad R_x(\pi)R_y(\pi/2)=-iH\quad\text{（01空间）}.
\end{equation}
异位单门可并行，不相交CZ可并行，但单门层与CZ层隔离。虚拟参考帧不占脉冲时间，更新顺序仍决定后续轴。
\datafig{00_native_primitives}{CNOT、H和完整01空间SWAP的原生展开。X是原生$X_{01}$，没有拆成两个半门；CZ列与单门列隔离。}{fig:primitives}

\subsection{节点只需在规定入口上正确}
下行只要求$L=R=0$时
\begin{equation}\label{eq:router}
 U_{\downarrow}\ket{c,i,0,0}=
 \begin{cases}\ket{0,0,i,0},&c=0,\\\ket{1,0,0,i},&c=1.\end{cases}
\end{equation}
四列合法入口与四行目标像之间，理想矩阵是
\begin{equation}
\begin{array}{c|cccc}
 &0000&0100&1000&1100\\\hline
0000&1&0&0&0\\0010&0&1&0&0\\1000&0&0&1&0\\1001&0&0&0&1
\end{array}.
\end{equation}
这是四维子空间的等距映射，不是对所有16个计算基入口定义的通用受控交换。上行需在该像空间上恢复原入口及相位。整树能用它，前提是每次调用始终满足入口约束。
\FloatBarrier

\section{R5与R6怎样处理临时2态}\label{sec:routers}
\subsection{R5：编码后分别搬运}
先定义
\begin{equation}\label{eq:encoder}
 E:X_{01}(I)\seq X_{12}(I)\seq C_{C\to I}\seq X_{12}(I)\seq X_{01}(I).
\end{equation}
R5下行按时间执行
\begin{equation}\label{eq:r5down}
\begin{aligned}
 &E\seq V_{\downarrow}(I)\seq T_{I\to L}\seq V_{\downarrow}^{\dagger}(L)\\
 &\seq X_{12}(I)\seq T_{I\to R}\seq F_{\rm out}.
\end{aligned}
\end{equation}
$V$的目标012作用是$V\oplus1$，允许公共全局相位。它由两个01半门和两跃迁帧展开，不增加CZ。编码用1个CNOT，两次搬运各2个，合计5个CZ。

暂取理想CZ的高能级旁路为恒等、$V=I$，保留$\pi$门相位，路径如下：
\begin{center}\small\begin{tabular}{@{}cccc@{}}\toprule
入口CI & 编码后CI & 左搬运后CILR & 右搬运后CILR\\\midrule
$00$ & $+\ket{00}$ & $+\ket{0000}$ & $+\ket{0000}$\\
$01$ & $-\ket{01}$ & $-\ket{0010}$ & $-\ket{0010}$\\
$10$ & $+\ket{10}$ & $+\ket{1000}$ & $+\ket{1000}$\\
$11$ & $-\ket{12}$ & $-\ket{1200}$ & $+i\ket{1001}$\\\bottomrule
\end{tabular}\end{center}
$C=0,I=1$先送左；$C=1,I=1$暂存于2态，经过左搬运后由$X_{12}$变回1，再送右。理想端点帧$(\phi_C,\phi_L,\phi_R)=(0,-\pi,\pi/2)$消去相对相位。这是理想矩阵解，不是实际脉冲校准值。
\datafig{03_native_r5}{R5下行原生线路。两行连续阅读，没有重置；蓝块为01门，红块为12门。实际轴和参考帧见附录逐门表。}{fig:r5}

\subsection{2态作用在20/21相位，而不是必然破坏11校准}
01空间校准主要固定
\begin{equation}
 \chi_{11}-\chi_{10}-\chi_{01}+\chi_{00}\simeq\pi,
\end{equation}
却没有同时固定$\chi_{21}-\chi_{20}$。第一位指临时含2态的线；若物理位顺序相反，应改写为12/02扇区。

含2态的控制经过CNOT目标半门夹层，看到的目标块近似为
\begin{equation}\label{eq:B}
 B=R_y(\pi/2)\begin{pmatrix}e^{i\chi_{20}}&0\\0&e^{i\chi_{21}}\end{pmatrix}R_y(-\pi/2).
\end{equation}
令$\bar\chi=(\chi_{20}+\chi_{21})/2$、$\delta=\chi_{21}-\chi_{20}$，则
\begin{equation}\label{eq:Baction}
 B\ket0=e^{i\bar\chi}\left[\cos(\delta/2)\ket0-i\sin(\delta/2)\ket1\right].
\end{equation}
旁路不再静止，而可能改变目标人口。\textbf{末端对角虚拟Z只能改相位，不能消掉已经产生的错误人口。}这与“11相位校准失败”是两件事。真实CZ还可能发生$\ket{21}$与$\ket{30}$等振幅交换，完整脉冲保留这些作用；式\eqref{eq:B}只是近似封闭块的解释。

\subsection{配对V怎样同时保留正常搬运、恢复辅助路径}
令左搬运$K=T_{I\to L}$。正常01入口近似满足$K\ket{\psi,0}\simeq\ket{0,\psi}$；辅助路径近似为$K\ket{2,0}\simeq\ket2 M\ket0$。若该块近似封闭、泄露小，选择V的第一列匹配$M\ket0$的归一化方向，则
\begin{align}
 \ket{\psi,0}&\xrightarrow{V(I)}\ket{V\psi,0}
 \xrightarrow{K}\ket{0,V\psi}\xrightarrow{V^\dagger(L)}\ket{0,\psi},\label{eq:pairnormal}\\
 \ket{2,0}&\xrightarrow{V(I)}\ket{2,0}
 \xrightarrow{K}\ket2M\ket0\xrightarrow{V^\dagger(L)}e^{i\gamma}\ket{2,0}.\label{eq:pairaux}
\end{align}
第一行正常搬运保留；第二行辅助目标返回空态，剩余分支相位再由端点校准处理。漏出该块的振幅无法由V恢复，结论也依赖$L$初始0及正常搬运接近理想。

这是利用受限入口的\textbf{配对匹配}，没有新增回波CZ，不是通用完整CSWAP恒等式。仅当人口路径已经正确、剩余相位可分解为局域相位时，输出虚拟Z才足够。

\subsection{R6和实际回程模板}
R6下行使用不同重排，不含配对V：
\begin{equation}\label{eq:r6}
\begin{aligned}
 &C_{I\to L}\seq X_{01}(I)\seq X_{12}(I)\seq C_{C\to I}\\
 &\seq X_{12}(I)\seq X_{01}(I)\seq C_{L\to I}\seq X_{12}(I)\\
 &\seq C_{I\to L}\seq C_{I\to R}\seq C_{R\to I}\seq F_{\rm out}.
\end{aligned}
\end{equation}
每个方向6个CZ、17个物理单门；R5为5个CZ、19个单门。不用V不代表高能级物理作用消失，仍需脉冲和往返验证。
\datafig{03_native_r6}{R6下行原生线路。以一个额外CZ换取不使用配对V；单门与CZ仍分层隔离。}{fig:r6}

R5上行按时间执行
\begin{equation}\label{eq:r5up}
\begin{aligned}
 &F_{\rm in}\seq T_{R\to I}\seq X_{12}(I)\seq V_{\uparrow}^{\dagger}(L)\\
 &\seq T_{L\to I}\seq V_{\uparrow}(I)\seq\operatorname{rev}_{\CZ}(E).
\end{aligned}
\end{equation}
R6上行先加输入端点帧，再执行裸下行的$\operatorname{rev}_{\CZ}$。该操作倒转顺序、将01半门轴加$\pi$、显式帧取反，但$\pi$门与CZ仍用同一物理脉冲。

它\textbf{不是完整多能级传播子的数学dagger}：高能级CZ不保证自逆，$\pi$门也有旁路相位。因此上下行分别保存校准参数，不能默认$V_{\uparrow}=V_{\downarrow}$，还必须直接验证实际down接实际up。端点帧下行放输出，上行放输入。
\FloatBarrier

\section{从三个节点到固定数据库查询}\label{sec:query}
\subsection{17个角色和共享端口}
QPU有$A_0,A_1,D$。临时树有三个控制$C_0,C_1,C_2$；数据端口包括根与子输入$T_0,T_1,T_2$，以及四个叶端$L_{00},L_{01},L_{10},L_{11}$。另有四个存储$M_{00},M_{01},M_{10},M_{11}$。节点映射为
\begin{center}\small\begin{tabular}{@{}ccccc@{}}\toprule
节点&C&I&L&R\\\midrule
根&$C_0$&$T_0$&$T_1$&$T_2$\\左子&$C_1$&$T_1$&$L_{00}$&$L_{01}$\\右子&$C_2$&$T_2$&$L_{10}$&$L_{11}$\\\bottomrule
\end{tabular}\end{center}
根输出与子节点输入共用物理端口；$T_1,T_2$在这里是端口名，与寿命参数无关。左右子节点互不重叠可并行，根与共享端口的子节点不能同时占用。
\datafig{01_architecture}{17角色和参考连边。$D-T_0$是唯一计入的跨片接口；还假定QPU具有$A_0-D$、$A_1-D$及$A_0-A_1$边，叶端各有存储访问边。}{fig:architecture}

\subsection{网关如何恢复中间位D}
定义
\begin{equation}
 G(A_j):S_{A_jD}\seq S_{DT_0}\seq S_{A_jD}.
\end{equation}
初始标签$(a,d,t)$经历$(a,d,t)\to(d,a,t)\to(d,t,a)\to(t,d,a)$，所以交换$A_j$与$T_0$，并恢复D。线性扩展后对01叠加及纠缠也成立。每次G用9个CZ，其中3个跨片。

\subsection{完整顺序：加载、访问、恢复}
表\ref{tab:querysequence}按实际时间列出宏观线路。$\parallel$表示并行，其余串行。临时树入口全0；每次下行的目标端口由前面的保存与恢复步骤保持为空。
{\small\setlength{\tabcolsep}{4pt}
\begin{longtable}{@{}r p{.44\linewidth} p{.43\linewidth}@{}}
\caption{完整query的宏观时间顺序}\label{tab:querysequence}\\\toprule
步 & 操作（时间由上到下） & 目的与约束\\\midrule\endfirsthead
\multicolumn{3}{c}{续表：完整query的宏观时间顺序}\\\toprule
步 & 操作（时间由上到下） & 目的与约束\\\midrule\endhead
\midrule\multicolumn{3}{r}{下页续}\endfoot
\bottomrule\endlastfoot
1 & $G(A_0)$ & 地址高位到根输入，D恢复\\
2 & $S_{T_0C_0}$ & 高位存入根控制，根输入回0\\
3 & $G(A_1)$ & 低位到根输入，D恢复\\
4 & 根router down & 按根控制把低位送至左右子输入\\
5 & $S_{T_1C_1}\parallel S_{T_2C_2}$ & 保存低位，两输入端口回0\\
6 & $H_D=Y_{01}^{(D)}/2\to X_{01}^{(D)}$ & 将异或读取转为相位形式\\
7 & $S_{DT_0}$ & 数据进入根；跨片\\
8 & 根router down & 选择左右子树\\
9 & 左down $\parallel$ 右down & 数据至对应叶端\\
10 & $\prod_j\CZ(L_j,M_j)$（四门并行） & 固定访问：只有被寻址叶端携带数据\\
11 & 左up $\parallel$ 右up & 数据返至子输入，叶端恢复\\
12 & 根router up & 数据返回根输入\\
13 & $S_{DT_0}$ & 数据返回QPU；跨片\\
14 & $S_{T_1C_1}\parallel S_{T_2C_2}$ & 取出低位，子控制恢复\\
15 & 根router up & 低位返回根输入\\
16 & $G(A_1)$ & 恢复A1\\
17 & $S_{T_0C_0}$ & 取出高位，根控制恢复\\
18 & $G(A_0)$ & 恢复A0，临时树回零\\
19 & $H_D=Y_{01}^{(D)}/2\to X_{01}^{(D)}$ & 解码为异或query\\
\end{longtable}}

共8次router调用：根用于地址低位加载/卸载2次，加数据经过三个节点的下行/上行6次；合计下行4次、上行4次。不能只计数据路径而漏掉地址处理。
\datafig{02_pipeline}{一次完整query的因果顺序。地址保存与卸载、数据下行与返回缺一不可；上行采用实际独立参数，而不是数值强行取下行dagger。}{fig:pipeline}

\subsection{四个固定存储CZ为何只访问一个地址}
把理想地址保存和数据路由编码记为W。H编码后的数据基态记为x，到达访问层时有$n_{L_j}=x\delta_{j,a}$，因此
\begin{equation}\label{eq:memoryphase}
 P_M=\prod_j\CZ(L_j,M_j):\quad
 \ket{a,x}\mapsto(-1)^{\sum_jn_{L_j}m_j}\ket{a,x}=(-1)^{xm_a}\ket{a,x}.
\end{equation}
返回后临时树清零，两侧H将相位作用转为异或：
\begin{equation}\label{eq:queryconstruction}
 Q_m=H_DW^\dagger P_MWH_D,\qquad Q_m\ket{a,d}=\ket{a,d\oplus m_a}.
\end{equation}
该式用于理想推导。实际表中入口H在地址加载之后；理想加载恢复D，故与$H_D$可交换。脉冲实现仍按表传播，实际up不被数学dagger替代。

例如$M=0010$只翻转地址10的数据位。按$\ket{A_0A_1D}$基序，
\begin{equation}\label{eq:matrixexample}
 Q_{0010}=\diag(I_2,I_2,X,I_2)=
 \begin{pmatrix}
1&0&0&0&0&0&0&0\\0&1&0&0&0&0&0&0\\0&0&1&0&0&0&0&0\\0&0&0&1&0&0&0&0\\
0&0&0&0&0&1&0&0\\0&0&0&0&1&0&0&0\\0&0&0&0&0&0&1&0\\0&0&0&0&0&0&0&1
\end{pmatrix}.
\end{equation}
数据库切换只改M初态，访问层四个CZ和主体线路都不改。模拟将基态存储精确收缩为叶端Z或I，不代表硬件计划删除存储门。

\subsection{先证明理想结构}
检查逐项相乘原生单门、CZ和参考帧，没有把节点当未展开黑箱。R5/R6各16数据库的完整$8\times8$复数query矩阵均通过，只去掉一个公共全局相位，最大元素误差为$\IdealError$。

另外对128个地址/数据/存储组合与256个H编码分支逐步核验，共2048次router调用均满足合法入口，最后辅助全零、数据库保持。由此才可将部分节点用于完整查询，而不是靠某个特定输入偶然通过。
\FloatBarrier

\section{从理想门到冻结脉冲：本次模拟了什么}\label{sec:pulse}
\subsection{局域参考器件与波形}
每个局域节点用四个Duffing振子，每位保留0、1、2、3四能级。频率和非谐性如下，来自合成参考器件，不是最新芯片校准。
\begin{table}[htbp]\centering\small
\caption{局域四能级参考器件}\label{tab:device}
\begin{tabular}{@{}llrr@{}}\toprule
线 & 参考标签 & $f_{01}$/GHz & $\alpha$/GHz\\\midrule
C & q111 & 4.600 & -0.220\\
I & q121 & 5.000 & -0.220\\
L & q130 & 5.500 & -0.220\\
R & q112 & 5.650 & -0.235\\
\bottomrule\end{tabular}\end{table}

在公共旋转系中，Hamiltonian结构为
\begin{equation}\label{eq:hamiltonian}
\begin{aligned}
 H/\hbar={}&2\pi\sum_q\left[(f_q-f_{\rm rot})n_q+\frac{\alpha_q}{2}n_q(n_q-1)\right]\\
 &+2\pi\sum_{(a,b)}g_{ab}(t)(a_a^\dagger a_b+a_a a_b^\dagger)
 +\pi\sum_q[\widetilde\Omega_q(t)a_q^\dagger+\widetilde\Omega_q^*(t)a_q].
\end{aligned}
\end{equation}
频率单位GHz，时间ns，Hamiltonian角频率rad/ns。$\widetilde\Omega$包含包络、轴相位和载波参考。单门时耦合关闭，CZ只打开选中边，其余边为0；这使不相交子块积分可以精确张量应用到四体态，而没有忽略选中边的高能级作用。

01门长32 ns，12门长40 ns。余弦包络与DRAG为
\begin{equation}
 s(t)=\frac{1-\cos(2\pi t/T)}2,\qquad
 \Omega(t)=A\left[s(t)-i\eta\frac{\dot s(t)}{\alpha}\right].
\end{equation}
$A,\eta$按模拟器约定定义，不是可直接上传的AWG电压。CZ用起、落半余弦权重w同时调频和开耦合：
\begin{equation}
 f_q(t)=f_q^0+w(t)(f_q^{\CZ}-f_q^0),\qquad
 g_{ab}(t)=\rho_{ab}\sqrt{f_a(t)f_b(t)}\,w(t)p.
\end{equation}
三边$\rho=0.0031$、$p=1$。当前三条CZ时长约30.642、56.371、64.820 ns，详细参数见附录\ref{app:parameters}。
\datafig{27_waveform}{R5下行微波包络、频率与耦合轨迹。图省略快速载波，不代表相位为零。积分步长0.1 ns，绘图采样约0.5 ns。}{fig:waveform}

\subsection{虚拟相位是参考帧更新}
每条线维护软件缓冲$(f_{01},f_{12})$。01门使用轴$\phi+f_{01}$，12门使用$\phi+f_{12}$。\code{frame(phi)}同时更新两者；\code{frame_01}、\code{frame_12}只更新相应跃迁。单门和CZ的原生补偿也自动累计到缓冲。

去掉已知空闲参考旋转后，对应的逻辑帧定义是
\begin{equation}\label{eq:frame}
 F(f_{01},f_{12})=\diag(1,e^{-if_{01}},e^{-i(f_{01}+f_{12})},e^{-i(f_{01}+2f_{12})}).
\end{equation}
普通帧$\phi$在012空间为$\diag(1,e^{-i\phi},e^{-2i\phi})$，不是任意指定两个能级相位。R5的V展开需要两跃迁缓冲。它们是后续驱动参考更新，不是瞬间物理Z踢；拼接模块不能重置帧，也不能重复加入已由门库实施的补偿。

\subsection{整树是局域脉冲导出的模块过程}
局域门先经四能级波形积分，再作为有效过程嵌入整树。模块外搬运、QPU、跨片与存储门仍为理想01过程延拓；四个经典基态存储作精确条件变量。量子传播显式保留13个四能级站点，主结果使用最大bond 128的张量网络。

因此后文数字应称为\textbf{局域脉冲导出的整机相干基线}，而非17体连续全波形结果。尚无一套满足全部共享端口频率的17位Hamiltonian，也没有把所有接口门换成真实脉冲。该边界限制了对外部泄露和真机精度的解释。
\FloatBarrier

\section{先检查单节点：人口、相干与实际往返}\label{sec:local}
\subsection{临时2态不等于最终泄露}
编码会主动占据2态。图\ref{fig:temporary}对真空入口显示原生层边界人口；$P_2$中途接近1可以是设计动作。应判断末态是否回到规定逻辑像空间，以及返回时相位是否正确，而不能把这段曲线直接称为100\%泄露错误。
\datafig{13_temporary_level2}{临时2态占据，仅采样原生层边界，不是CZ内部瞬时峰值。此模型没有耗散帮助复位。}{fig:temporary}

\subsection{连续控制与16态互相补充}
固定$I=1,L=R=0$，将控制从0连续变化到1：
\begin{equation}
 \ket{\psi_C(\theta)}=\cos(\theta/2)\ket0+\sin(\theta/2)\ket1.
\end{equation}
理想输出是$\cos(\theta/2)\ket{0010}+\sin(\theta/2)\ket{1001}$，因此左右端$P_1$分别为$\cos^2(\theta/2)$和$\sin^2(\theta/2)$。101个角度共用同一套冻结波形。
\datafig{12_control_sweep}{连续控制扫描的四位末态人口。平滑分流说明路径合理，相干性还需联合态与通道检查，不能单由人口推出。}{fig:sweep}

C与I各取$\{0,1,+,-\}$形成16初态。down比较CLR联合态；up从理想合法像返回CI；实际往返使用真实down输出接真实up。表\ref{tab:localresults}另外保留完整四线合法块的过程$\Fe$，不能与约化态重叠混为一谈。
\begin{table}[htbp]\centering\small
\caption{局域原始指标；16态重叠与合法四线过程是不同检验}\label{tab:localresults}
\begin{tabular}{@{}lrrr@{}}\toprule
方案/方向 & 16态平均重叠 & 合法过程$\Fe$ & 01生存率\\\midrule
R5 DOWN & 99.6244\% & 99.5699\% & 99.5732\%\\
R5 UP & 99.5742\% & 99.5596\% & 99.5729\%\\
R5 ROUNDTRIP & 99.3226\% & 99.3102\% & 99.3282\%\\
R6 DOWN & 99.6068\% & 99.5516\% & 99.5706\%\\
R6 UP & 99.5682\% & 99.5560\% & 99.5702\%\\
R6 ROUNDTRIP & 99.4441\% & 99.4171\% & 99.4515\%\\
\bottomrule\end{tabular}\end{table}

\datafig{10_local_truth}{局域四列合法基态的转移人口。up列为理想合法像；行列以逻辑CI标签标记。未按列归一，缺失可以是高能级，也可以是错误01态。}{fig:localtruth}
\datafig{11_local_states}{16态原始目标重叠：down、独立合法入口up及实际down接实际up。全部为直接传播，不含有限shots或读出模型。}{fig:localstates}

局域下行约99.6\%，实际往返略低，且两方案排序随检验方向改变。额外一个CZ不是唯一因素，回程、分支相位与相干干涉同样决定表现。原上行参数标定的相干目标是匹配down准备的往返；这里另外计算独立up，不把原标定说成独立up全面优化。

\subsection{配对修正的消融证据}
固定CZ和端点帧，移除R5的两个配对V，不重新优化其他参数。$CI=11$的完整四线目标人口由\RfivePairPop{}降为\RfiveNoPairPop{}，合法块$\Fe$由\RfiveLocalFe{}降为\RfiveNoPairFe{}。
\datafig{14_paired_ablation}{R5配对消融，与式\eqref{eq:pairnormal}、\eqref{eq:pairaux}对应。配对恢复当前辅助路径的目标旋转，没有新增CZ；该消融不证明所有无配对构造都不可能。}{fig:pairablation}
\FloatBarrier

\section{再检查完整query：先功能，后诊断}\label{sec:network}
\subsection{人口与含相位矩阵一起看}
对$M=0010$，理想矩阵已在式\eqref{eq:matrixexample}写明。图\ref{fig:truth8}迹掉所有临时树，展示输出人口；图\ref{fig:k0}展示真实辅助零分支的复数操作块
\begin{equation}\label{eq:K0}
 K_0=P_{AD,01}\bra0_{\mathcal R}U_m\ket0_{\mathcal R}P_{AD,01}.
\end{equation}
前者不要求辅助回零，后者保留返回分支的原始幅度与相位。二者来自同一次传播。$K_0$通常不酉：投影后概率去了其他Kraus块，不能将每列归一化来掩盖损失。
\datafig{24_query_truth}{计算基$8\times8$人口矩阵，行输出、列输入，迹掉全部辅助态。它接近目标仍不能证明相位正确。}{fig:truth8}
\datafig{23_query_matrix}{同一次query的$K_0$实、虚部，只去掉公共全局相位，不按列归一。实部色标约$\pm1$，虚部约$\pm0.08$，颜色深浅不能直接跨图比较量级。}{fig:k0}

\subsection{全部数据库与完整输入通道}
16数据库均用固定线路。图\ref{fig:querystates}先给9种地址入口、$D=0$的联合态重叠。为覆盖未逐一展示的所有复数输入，再使用三输入Choi态。
\datafig{20_query_states}{16数据库与9地址入口的原始联合态重叠，无后选择和计算块重归一。色标从0.9开始以显示相干残差。}{fig:querystates}
将三输入AD与数学参考制备最大纠缠，经query后迹掉树；取AD输出01块为J，定义
\begin{equation}\label{eq:Fe}
\begin{aligned}
 \ket{\Phi_8}&=\frac1{\sqrt8}\sum_{x=0}^7\ket{x,x},\qquad
 \ket{\Phi_{Q_m}}=(I\otimes Q_m)\ket{\Phi_8},\\
 \Fe&=\bra{\Phi_{Q_m}}J\ket{\Phi_{Q_m}}.
\end{aligned}
\end{equation}
参考只是数学工具，不计入硬件位数。R5/R6平均原始$\Fe$为\RfiveFe{}与\RsixFe{}，它比某个算法输入成功更全面。
\datafig{21_query_process}{完整query的原始$\Fe$、辅助全0概率与条件Choi重叠。条件值只用于诊断，不能替代原始确定性性能。}{fig:queryprocess}
辅助零Choi块记为$J_0$，接受率$P_0=\tr J_0$，条件Choi重叠为$\Fe^{(0)}/P_0$。应分开报告原始、接受率和条件值；逐数据库归一后取平均，也不能用几个均值相乘代替联合概率。

\subsection{剩余误差落在哪些分支}
令$S=\tr J$为AD计算空间生存率，$\Fe^{(0)}=\bra{\Phi_{Q_m}}J_0\ket{\Phi_{Q_m}}$，则
\begin{equation}\label{eq:errorpartition}
\begin{aligned}
1-\Fe={}&\underbrace{1-S}_{\text{AD计算空间外}}
+\underbrace{S-P_0-(\Fe-\Fe^{(0)})}_{\text{辅助非零分支内不匹配}}\\
&+\underbrace{P_0-\Fe^{(0)}}_{\text{辅助零分支内不匹配}}.
\end{aligned}
\end{equation}
这是非负通道分解（浮点误差除外），不是人为设定门错误概率。R5后两项约1.953与0.063个百分点，R6约1.541与0.151个百分点。主要残差落在辅助恢复，但不能据此认定某一个X12是源头；非零辅助也可能是错误01态。
\datafig{22_coherent_residual}{相干残差分解。AD泄露接近零部分来自模块外理想01延拓，不能解释成真机泄露已经解决。}{fig:budget}
\datafig{33_layer_tracking}{同一query的阶段末人口，数据库0010、地址$\ket{++}$、数据$\ket-$。中途控制激发与端口2态可属设计，关键是解码后恢复。它用于定位异常阶段，不直接确定物理机制。}{fig:tracking}

\subsection{末端虚拟Z能修到什么程度}
只在辅助0且D返回$\ket-$的分支，扣去理想数据库标记，定义四个地址残余相位$\epsilon_a$。两个末端地址局域Z不能改变
\begin{equation}
 \Delta=\epsilon_{00}-\epsilon_{01}-\epsilon_{10}+\epsilon_{11}\pmod{2\pi}.
\end{equation}
非零$\Delta$提示该返回分支有非局域相位。该诊断不能把原始混合输出硬解释成四个相位；末端对角帧也不能提高辅助回零概率。中间驱动轴可能改变返回干涉，需重新检查完整过程。
\datafig{25_phase_invariant}{返回分支的相位不变量：先投影辅助0、D负号态，再扣除理想标记。它对两个地址末端Z不变。}{fig:invariant}
\FloatBarrier

\section{最后闭合算法：从相干查询到Grover}\label{sec:algorithm}
\subsection{先用GHZ检验跨地址相干}
地址制备$(\ket{00}+\ket{11})/\sqrt2$、D为0、数据库0101，目标AD为
\begin{equation}
 \ket{\mathrm{GHZ}}=\frac{\ket{000}+\ket{111}}{\sqrt2}.
\end{equation}
除了两个角点人口，$\rho_{000,111}$也必须保留。原始目标重叠为R5的\RfiveGHZ{}与R6的\RsixGHZ{}。
\datafig{30_GHZ}{原始AD的GHZ密度矩阵。非对角角点检验相干，虚部显示相对相位。未后选择，也不含实验准备和分析误差。}{fig:ghz}

\subsection{相位回踢与扩散}
数据$\ket-=(\ket0-\ket1)/\sqrt2$满足$X\ket-=-\ket-$，所以同一query实现
\begin{equation}
 Q_m\ket{a,-}=(-1)^{m_a}\ket{a,-}.
\end{equation}
均匀地址$\ket s=\ket{++}$经相位oracle和$2\ket s\bra s-I$后，四选一理想搜索一轮即到目标。当前扩散的时间顺序是
\begin{equation}\label{eq:diffusion}
\begin{aligned}
 &(H_{A_0}\parallel H_{A_1})\seq(X_{A_0}\parallel X_{A_1})\seq\CZ_{A_0A_1}\\
 &\seq(X_{A_0}\parallel X_{A_1})\seq(H_{A_0}\parallel H_{A_1}).
\end{aligned}
\end{equation}
每个H按式\eqref{eq:h}展开；该序列与标准扩散只差公共负号。
\datafig{32_algorithm_stages}{R5目标10：入口、query后、扩散后。oracle主要改变目标相关相干符号，人口仍近似均匀；扩散再将符号差转为目标人口。来自同一次传播，计算块未另归一。}{fig:algorithmstages}

\subsection{四目标同线路验证}
四个任务只改数据库为1000、0100、0010、0001，query和扩散相同。平均原始地址成功率为\RfiveGrover{}与\RsixGrover{}。表\ref{tab:grover}同时保留辅助零接受率、条件成功与联合正确比例。
\begin{table}[htbp]\centering\small
\caption{四目标Grover：条件高值必须与接受率及联合比例一起读}\label{tab:grover}
\begin{tabular}{@{}lrrrr@{}}\toprule
方案/目标 & 原始成功 & 辅助0概率 & 条件成功 & 联合正确\\\midrule
R5/00 & 98.8625\% & 98.5175\% & 99.9548\% & 98.4730\%\\
R5/01 & 98.7905\% & 98.4163\% & 99.9545\% & 98.3715\%\\
R5/10 & 99.1900\% & 99.0228\% & 99.9535\% & 98.9768\%\\
R5/11 & 98.7596\% & 98.3955\% & 99.9556\% & 98.3519\%\\
R6/00 & 98.7949\% & 98.6548\% & 99.8518\% & 98.5085\%\\
R6/01 & 98.6297\% & 98.4050\% & 99.8465\% & 98.2539\%\\
R6/10 & 98.6653\% & 98.3670\% & 99.8540\% & 98.2234\%\\
R6/11 & 98.8048\% & 98.5733\% & 99.8553\% & 98.4307\%\\
\bottomrule\end{tabular}\end{table}

\datafig{31_Grover}{四目标原始Grover人口，行是标记地址，列是实际输出地址。没有shots抽样、读出分类或附加随机门事件，只支持该参考过程的单轮功能。}{fig:grover}
本次输入$\ket{++,-}$与数据库直接初始化。真实准备、复位和读出脉冲未评估，资源也应另加，不能将这里的成功率称为完整实验保真度。
\FloatBarrier

\section{需要多少资源，哪些地方值得压缩}\label{sec:resource}
\begin{table}[htbp]\centering\small
\caption{两层四存储的一轮Grover；不含真实准备、复位与读出}\label{tab:resources}
\begin{tabular}{@{}lrr@{}}\toprule
资源 & R5 & R6\\\midrule
总角色数 & 17 & 17\\
router调用 & 8 & 8\\
router内部CZ & 40 & 48\\
CZ总数 & 105 & 113\\
其中跨片CZ & 18 & 18\\
物理单门 & 288 & 272\\
CZ层深 & 86 & 92\\
总调度层（含入口等待） & 283 & 289\\
query时长/ns & 11184.152 & 11522.380\\
含扩散总时长/ns & 11426.152 & 11764.380\\
\bottomrule\end{tabular}\end{table}

R5/R6一轮分别105/113个CZ：节点部分是$8\times5/8\times6=40/48$，其余65个来自树控制保存与取出18、跨片完整SWAP 18、QPU中转SWAP 24、存储访问4、扩散1。跨片CZ已在总数中，不能重复加。

当前20个完整SWAP贡献60个CZ和120个半门。即使节点已经较短，接口与地址处理仍占很大比例。接收端已知0时可研究两CNOT搬运替换，但须重新证明全程入口、回程与相位，不能在每个位置直接删第三个CNOT。
\datafig{07_cost}{同一拓扑下的资源。10-CZ柱只是假定每次router核用10个CZ、保留其他协议的估算（$8\times10+65=145$），没有在此验证其脉冲性能，也不是最优门数证明。}{fig:cost}
对n层、$N=2^n$存储，本协议为
\begin{equation}
 N_{\rm role}=4N+n-1,\qquad N_{\rm router}=4N-2n-4.
\end{equation}
n为1、2、3时，角色数8/17/34，router调用2/8/22。这是串行地址注入、显式端口方案的实际计数，不是下界。
\datafig{08_resources_scale}{1至3层资源。时间统一比较query；CZ总数只有n=2含一次扩散，该点比较纯查询时应减1。H均已实际展开。}{fig:resourcescale}
\FloatBarrier

\section{数值可信度与结论边界}\label{sec:boundary}
\subsection{积分、能级与张量网络}
主局域步长0.1 ns，另固定参数用0.05 ns，并将四能级增加至五能级，没有重新优化。图\ref{fig:convergence}的变化小于百分比量级残差，支持局域传播与截断足够稳定。
\datafig{35_pulse_convergence}{四种上下行的积分与能级检查。五能级不是新的优化门，只是同一冻结波形的截断敏感性检验。}{fig:convergence}
整网使用bond 128，另对目标10用64重算；R5/R6查询$\Fe$差约$2.18\times10^{-6}$与$3.79\times10^{-5}$。同时检查正性、迹、辅助选择界与完整原生矩阵。累计截断权重只是诊断，不能直接当作保真度误差。

\subsection{扩大规模仍会累积相干残差}
固定均匀地址、D负号态，测试全1与单目标两种数据库模式。图\ref{fig:scaling}说明无耗散时模块相干偏差仍会累积。三层没有覆盖全部数据库、全部输入，也没有严格误差界。
\datafig{34_scaling}{R5的1至3层相位query。三层原始目标重叠约96.5\%--97.1\%，均用bond 128；只是两种数据库模式的一次query，不是多轮Grover或可扩展性证明。}{fig:scaling}

\subsection{真实器件还缺哪层证据}
最重要的缺口是共享工作点。根输出与子节点中心为同一物理位，不能同时赋予两个参考角色的不同频率。真实系统应按节点实际工作点分别校准，必要时加入转场、等待相位和帧转换。

QPU搬运、跨片、存储访问及H展开尚未全部换成真实脉冲；外部AD几乎无泄露受理想01延拓限制。当前只研究基态经典存储，没有证明任意量子数据库读写。

这里只验证一轮，不能将成功率简单乘方预测重复查询。小的末态辅助残留会改变下一次部分router入口，必须直接验证多次复用。
\FloatBarrier

\section{综合判断与实验推进顺序}\label{sec:conclusion}
线路结构上，R5/R6在规定入口可构造正确路由；完整协议保持这些入口条件，固定访问与相干恢复给出query，四目标Grover闭环成立。

参考控制上，局域四能级与整树模块过程接近目标，但仍有辅助恢复和相位残差。V配对处理当前辅助路径的目标旋转，端点帧处理剩余分支相位；两者不能替代泄露与路径错误的修复。

建议按因果顺序推进：先确认17角色拓扑、共享端口与各节点工作点；再校准01/12门、CZ及含2态搬运块；随后做各节点down、独立up和实际往返；再做空query、路径Ramsey与辅助恢复；最后固定数据库query、GHZ和Grover。

这个顺序先排人口和回程，再排相位，最后进入算法，使失败有较小定位范围。当前结论是具备值得落地验证的线路，而非真实17位器件已被证明达到本文百分数。

\clearpage\appendix
\section{指标定义：原始、条件与联合量}\label{app:metrics}
理想目标纯态$\ket\psi$的原始态重叠是$F=\bra\psi\rho\ket\psi$。局域down比较CLR约化态，up及往返比较CI；合法过程块保留完整四线目标像。若四维块K按目标/入口态排序，则
\begin{equation}
 \Fe^{\rm legal}=\frac{|\tr K|^2}{16},\qquad
 P_{\rm image}=\frac14\tr(K^\dagger K).
\end{equation}
前者也检查相位，后者只看目标像人口，均未按列归一。

完整query的全输出Choi态是$J_{\rm full}=(I\otimes\mathcal E)(\proj{\Phi_8})$；输出限制AD的01块后为J，迹$S\leq1$，泄露时它不归一。Haar平均原始目标态重叠为
\begin{equation}
 \overline F_{\rm Haar}=\frac{8\Fe+S}{9}.
\end{equation}
只有S为1才退化成常用迹保持形式。

给定输入，真实辅助零投影$P_{\mathcal R0}$定义
\begin{equation}
 p_0=\tr(P_{\mathcal R0}\rho),\qquad
 \rho_0=P_{\mathcal R0}\rho P_{\mathcal R0}/p_0.
\end{equation}
目标地址测量为$\Pi_a$时，三种成功指标分别为
\begin{align}
 p_{\rm raw}&=\tr(\Pi_a\rho),\\
 p_{\rm cond}&=\tr(\Pi_a\rho_0),\\
 p_{\rm joint}&=\tr(\Pi_aP_{\mathcal R0}\rho P_{\mathcal R0}).
\end{align}
Choi平均接受率与某算法输入接受率不同，不能混用。

\clearpage\section{冻结参考参数与原生单门表}\label{app:parameters}
以下来自已有模拟JSON，计算用全精度，表中四舍五入。参考标签q111/q121/q130/q112中q121是中心I；不是完整17角色硬件映射，也不是当前硬件上传值。

附录末尾的四份局域门表中，$\phi_{\rm nom}$为协议相位；$\phi_{\rm buf}$为本地帧从零开始、累计补偿后的驱动轴，折回$[-\pi,\pi]$。它不含绝对载波，也不能直接当全局AWG相位。F类更新时长0；同层异位单门并行，层间隔离。
\subsection{R5的配对矩阵与实现}
\begin{equation}V_{\mathrm{down}}=\begin{pmatrix}0.923048+0.000000i & -0.001145-0.384684i\\0.274545-0.269459i & 0.648521+0.656839i\end{pmatrix}.\end{equation}
\begin{equation}V_{\mathrm{up}}=\begin{pmatrix}0.922191+0.000000i & 0.273149-0.273776i\\0.000251-0.386734i & 0.652412+0.651763i\end{pmatrix}.\end{equation}

对$V=e^{i\alpha}R_z(\beta)R_y(\theta)R_z(\gamma)$定义
\begin{equation}
 \widetilde Z(\delta):F_{01}(-\delta)\seq F_{12}(\delta/2).
\end{equation}
在012空间它等于嵌入$R_z(\delta)$乘公共相位。V的时间展开是
\begin{equation}
\begin{aligned}
 &\widetilde Z(\gamma)\seq\widetilde Z(\pi/2)\seq Y_{01}/2\seq\widetilde Z(\theta)
 \seq(-Y_{01}/2)\\
 &\seq\widetilde Z(-\pi/2)\seq\widetilde Z(\beta)\seq F_{12}(\alpha).
\end{aligned}
\end{equation}
末尾$F_{12}(\alpha)$维持01块相对2态的相位，不能丢弃。$V^\dagger$按矩阵重新分解；更高能级作用由原生展开与脉冲决定，不另假设恒等。

\begin{table}[htbp]\centering\small
\caption{单门冻结参数（程序幅度单位换算为MHz）}\label{tab:singleparams}
\begin{tabular}{@{}llrrrr@{}}\toprule
线 & 门 & T/ns & A/MHz & $\delta$/MHz & DRAG\\\midrule
C & $X_{01}/2$ & 32 & 15.623205 & -0.190561 & 0.144406\\
C & $X_{01}$ & 32 & 31.195108 & -1.076721 & 0.134111\\
C & $X_{12}$ & 40 & 17.692091 & -0.538408 & 0.082153\\
I & $X_{01}/2$ & 32 & 15.623236 & -0.190544 & 0.144409\\
I & $X_{01}$ & 32 & 31.156308 & -1.373690 & 0.149271\\
I & $X_{12}$ & 40 & 17.692086 & -0.538369 & 0.082149\\
L & $X_{01}/2$ & 32 & 15.623216 & -0.190535 & 0.144403\\
L & $X_{01}$ & 32 & 31.257109 & -0.507768 & 0.105243\\
L & $X_{12}$ & 40 & 17.692085 & -0.538399 & 0.082151\\
R & $X_{01}/2$ & 32 & 15.623367 & -0.180091 & 0.145040\\
R & $X_{01}$ & 32 & 31.271875 & -0.296084 & 0.095514\\
R & $X_{12}$ & 40 & 17.688586 & -0.545110 & 0.085622\\
\bottomrule\end{tabular}\end{table}
\begin{table}[htbp]\centering\small
\caption{CZ冻结参数，三边峰值因子均为1}\label{tab:czparams}
\begin{tabular}{@{}llrrr@{}}\toprule
边 & 调谐线 & $f^{\CZ}$/GHz & T/ns & 坡/ns\\\midrule
C--I & I & 4.818937693 & 30.641834 & 5.725930\\
I--L & L & 5.244278591 & 56.371311 & 11.894511\\
I--R & R & 5.259764008 & 64.820457 & 16.500000\\
\bottomrule\end{tabular}\end{table}


\begin{table}[htbp]\centering\small
\caption{端点帧：down放输出，up放输入，角度单位rad}\label{tab:endpoint}
\begin{tabular}{@{}lrrr@{}}\toprule
方案/方向 & $\phi_C$ & $\phi_L$ & $\phi_R$\\\midrule
R5 down & 2.136815 & -0.134507 & 2.348115\\
R5 up & 2.159159 & -0.153203 & 2.334839\\
R6 down & 2.139030 & -0.133450 & 2.956958\\
R6 up & 2.114272 & -0.143818 & 3.006746\\
\bottomrule\end{tabular}\end{table}


\clearpage{\small\setlength{\tabcolsep}{4pt}\fontsize{9}{10.6}\selectfont\renewcommand{\arraystretch}{1.0}
\begin{longtable}{@{}rrrllrrr@{}}
\caption{R5 down：18层，713.025368 ns}\label{tab:r5down}\\\toprule
序 & 层 & 起点/ns & 门 & 线 & $\phi_{\rm nom}$ & $\phi_{\rm buf}$ & 门长/ns\\\midrule\endfirsthead
\multicolumn{8}{c}{续表：R5 down：18层，713.025368 ns}\\\toprule
序 & 层 & 起点/ns & 门 & 线 & $\phi_{\rm nom}$ & $\phi_{\rm buf}$ & 门长/ns\\\midrule\endhead
\midrule\multicolumn{8}{r}{下页续}\endfoot
\bottomrule\endlastfoot
1 & 1 & 0.000 & $X_{01}$ & I & 0.00000 & 0.00000 & 32.000\\
2 & 2 & 32.000 & $X_{12}$ & I & 0.00000 & -0.01213 & 40.000\\
3 & 3 & 72.000 & $X_{01}/2$ & I & -1.57080 & -1.46338 & 32.000\\
4 & 4 & 104.000 & $\CZ$ & C,I & --- & --- & 30.642\\
5 & 5 & 134.642 & $X_{01}/2$ & I & 1.57080 & -1.49542 & 32.000\\
6 & 6 & 166.642 & $X_{12}$ & I & 0.00000 & -3.06430 & 40.000\\
7 & 7 & 206.642 & $X_{01}$ & I & 0.00000 & 3.08650 & 32.000\\
8 & 7 & 206.642 & $F_{01}$ & I & -1.56782 & --- & 0.000\\
9 & 7 & 206.642 & $F_{12}$ & I & 0.78391 & --- & 0.000\\
10 & 7 & 206.642 & $F_{01}$ & I & -1.57080 & --- & 0.000\\
11 & 7 & 206.642 & $F_{12}$ & I & 0.78540 & --- & 0.000\\
12 & 8 & 238.642 & $X_{01}/2$ & I & 1.57080 & 1.79488 & 32.000\\
13 & 8 & 238.642 & $F_{01}$ & I & -0.78973 & --- & 0.000\\
14 & 8 & 238.642 & $F_{12}$ & I & 0.39487 & --- & 0.000\\
15 & 9 & 270.642 & $X_{01}/2$ & I & -1.57080 & -2.09813 & 32.000\\
16 & 9 & 270.642 & $F_{01}$ & I & 1.57080 & --- & 0.000\\
17 & 9 & 270.642 & $F_{12}$ & I & -0.78540 & --- & 0.000\\
18 & 9 & 270.642 & $F_{01}$ & I & 0.77605 & --- & 0.000\\
19 & 9 & 270.642 & $F_{12}$ & I & -0.38802 & --- & 0.000\\
20 & 9 & 270.642 & $F_{12}$ & I & 0.39589 & --- & 0.000\\
21 & 9 & 270.642 & $X_{01}/2$ & L & -1.57080 & -1.57080 & 32.000\\
22 & 10 & 302.642 & $\CZ$ & I,L & --- & --- & 56.371\\
23 & 11 & 359.013 & $X_{01}/2$ & L & 1.57080 & -2.58760 & 32.000\\
24 & 11 & 359.013 & $X_{01}/2$ & I & -1.57080 & 0.54838 & 32.000\\
25 & 12 & 391.013 & $\CZ$ & L,I & --- & --- & 56.371\\
26 & 13 & 447.384 & $X_{01}/2$ & I & 1.57080 & -2.29354 & 32.000\\
27 & 13 & 447.384 & $F_{01}$ & L & 2.36554 & --- & 0.000\\
28 & 13 & 447.384 & $F_{12}$ & L & -1.18277 & --- & 0.000\\
29 & 13 & 447.384 & $F_{01}$ & L & -1.57080 & --- & 0.000\\
30 & 13 & 447.384 & $F_{12}$ & L & 0.78540 & --- & 0.000\\
31 & 13 & 447.384 & $X_{01}/2$ & L & 1.57080 & 0.33193 & 32.000\\
32 & 13 & 447.384 & $F_{01}$ & L & -0.78973 & --- & 0.000\\
33 & 13 & 447.384 & $F_{12}$ & L & 0.39487 & --- & 0.000\\
34 & 14 & 479.384 & $X_{01}/2$ & L & -1.57080 & 2.72210 & 32.000\\
35 & 14 & 479.384 & $F_{01}$ & L & 1.57080 & --- & 0.000\\
36 & 14 & 479.384 & $F_{12}$ & L & -0.78540 & --- & 0.000\\
37 & 14 & 479.384 & $F_{01}$ & L & -1.57377 & --- & 0.000\\
38 & 14 & 479.384 & $F_{12}$ & L & 0.78689 & --- & 0.000\\
39 & 14 & 479.384 & $F_{12}$ & L & -0.39589 & --- & 0.000\\
40 & 14 & 479.384 & $X_{12}$ & I & 0.00000 & -1.18282 & 40.000\\
41 & 14 & 479.384 & $X_{01}/2$ & R & -1.57080 & -1.57080 & 32.000\\
42 & 15 & 519.384 & $\CZ$ & I,R & --- & --- & 64.820\\
43 & 16 & 584.205 & $X_{01}/2$ & R & 1.57080 & 0.44708 & 32.000\\
44 & 16 & 584.205 & $X_{01}/2$ & I & -1.57080 & 0.97514 & 32.000\\
45 & 17 & 616.205 & $\CZ$ & R,I & --- & --- & 64.820\\
46 & 18 & 681.025 & $X_{01}/2$ & I & 1.57080 & -1.87059 & 32.000\\
47 & 18 & 681.025 & $F$ & C & 2.13681 & --- & 0.000\\
48 & 18 & 681.025 & $F$ & L & -0.13451 & --- & 0.000\\
49 & 18 & 681.025 & $F$ & R & 2.34811 & --- & 0.000\\
\end{longtable}}

\clearpage{\small\setlength{\tabcolsep}{4pt}\fontsize{9}{10.6}\selectfont\renewcommand{\arraystretch}{1.0}
\begin{longtable}{@{}rrrllrrr@{}}
\caption{R5 up：18层，713.025368 ns}\label{tab:r5up}\\\toprule
序 & 层 & 起点/ns & 门 & 线 & $\phi_{\rm nom}$ & $\phi_{\rm buf}$ & 门长/ns\\\midrule\endfirsthead
\multicolumn{8}{c}{续表：R5 up：18层，713.025368 ns}\\\toprule
序 & 层 & 起点/ns & 门 & 线 & $\phi_{\rm nom}$ & $\phi_{\rm buf}$ & 门长/ns\\\midrule\endhead
\midrule\multicolumn{8}{r}{下页续}\endfoot
\bottomrule\endlastfoot
1 & 1 & 0.000 & $F$ & C & 2.15916 & --- & 0.000\\
2 & 1 & 0.000 & $F$ & L & -0.15320 & --- & 0.000\\
3 & 1 & 0.000 & $F$ & R & 2.33484 & --- & 0.000\\
4 & 1 & 0.000 & $X_{01}/2$ & I & -1.57080 & -1.57080 & 32.000\\
5 & 2 & 32.000 & $\CZ$ & R,I & --- & --- & 64.820\\
6 & 3 & 96.820 & $X_{01}/2$ & I & 1.57080 & 1.86666 & 32.000\\
7 & 3 & 96.820 & $X_{01}/2$ & R & -1.57080 & -0.39589 & 32.000\\
8 & 4 & 128.820 & $\CZ$ & I,R & --- & --- & 64.820\\
9 & 5 & 193.641 & $X_{01}/2$ & R & 1.57080 & 1.62199 & 32.000\\
10 & 5 & 193.641 & $X_{12}$ & I & 0.00000 & 0.53957 & 40.000\\
11 & 5 & 193.641 & $F_{01}$ & L & 1.57144 & --- & 0.000\\
12 & 5 & 193.641 & $F_{12}$ & L & -0.78572 & --- & 0.000\\
13 & 5 & 193.641 & $F_{01}$ & L & -1.57080 & --- & 0.000\\
14 & 5 & 193.641 & $F_{12}$ & L & 0.78540 & --- & 0.000\\
15 & 5 & 193.641 & $X_{01}/2$ & L & 1.57080 & 1.41824 & 32.000\\
16 & 5 & 193.641 & $F_{01}$ & L & -0.79418 & --- & 0.000\\
17 & 5 & 193.641 & $F_{12}$ & L & 0.39709 & --- & 0.000\\
18 & 6 & 233.641 & $X_{01}/2$ & L & -1.57080 & -2.47922 & 32.000\\
19 & 6 & 233.641 & $F_{01}$ & L & 1.57080 & --- & 0.000\\
20 & 6 & 233.641 & $F_{12}$ & L & -0.78540 & --- & 0.000\\
21 & 6 & 233.641 & $F_{01}$ & L & -0.78654 & --- & 0.000\\
22 & 6 & 233.641 & $F_{12}$ & L & 0.39327 & --- & 0.000\\
23 & 6 & 233.641 & $F_{12}$ & L & -0.39245 & --- & 0.000\\
24 & 6 & 233.641 & $X_{01}/2$ & I & -1.57080 & -1.14785 & 32.000\\
25 & 7 & 265.641 & $\CZ$ & L,I & --- & --- & 56.371\\
26 & 8 & 322.012 & $X_{01}/2$ & I & 1.57080 & 2.29341 & 32.000\\
27 & 8 & 322.012 & $X_{01}/2$ & L & -1.57080 & 0.42982 & 32.000\\
28 & 9 & 354.012 & $\CZ$ & I,L & --- & --- & 56.371\\
29 & 10 & 410.384 & $X_{01}/2$ & L & 1.57080 & -0.58699 & 32.000\\
30 & 10 & 410.384 & $F_{01}$ & I & -2.35505 & --- & 0.000\\
31 & 10 & 410.384 & $F_{12}$ & I & 1.17752 & --- & 0.000\\
32 & 10 & 410.384 & $F_{01}$ & I & -1.57080 & --- & 0.000\\
33 & 10 & 410.384 & $F_{12}$ & I & 0.78540 & --- & 0.000\\
34 & 10 & 410.384 & $X_{01}/2$ & I & 1.57080 & -1.33276 & 32.000\\
35 & 10 & 410.384 & $F_{01}$ & I & -0.79418 & --- & 0.000\\
36 & 10 & 410.384 & $F_{12}$ & I & 0.39709 & --- & 0.000\\
37 & 11 & 442.384 & $X_{01}/2$ & I & -1.57080 & 1.05297 & 32.000\\
38 & 11 & 442.384 & $F_{01}$ & I & 1.57080 & --- & 0.000\\
39 & 11 & 442.384 & $F_{12}$ & I & -0.78540 & --- & 0.000\\
40 & 11 & 442.384 & $F_{01}$ & I & 1.57015 & --- & 0.000\\
41 & 11 & 442.384 & $F_{12}$ & I & -0.78507 & --- & 0.000\\
42 & 11 & 442.384 & $F_{12}$ & I & 0.39245 & --- & 0.000\\
43 & 12 & 474.384 & $X_{01}$ & I & 0.00000 & -0.48016 & 32.000\\
44 & 13 & 506.384 & $X_{12}$ & I & 0.00000 & 2.41640 & 40.000\\
45 & 14 & 546.384 & $X_{01}/2$ & I & -1.57080 & -1.94354 & 32.000\\
46 & 15 & 578.384 & $\CZ$ & C,I & --- & --- & 30.642\\
47 & 16 & 609.025 & $X_{01}/2$ & I & 1.57080 & -1.97559 & 32.000\\
48 & 17 & 641.025 & $X_{12}$ & I & 0.00000 & -0.63577 & 40.000\\
49 & 18 & 681.025 & $X_{01}$ & I & 0.00000 & 2.60634 & 32.000\\
\end{longtable}}

\clearpage{\small\setlength{\tabcolsep}{4pt}\fontsize{9}{10.6}\selectfont\renewcommand{\arraystretch}{1.0}
\begin{longtable}{@{}rrrllrrr@{}}
\caption{R6 down：19层，769.396679 ns}\label{tab:r6down}\\\toprule
序 & 层 & 起点/ns & 门 & 线 & $\phi_{\rm nom}$ & $\phi_{\rm buf}$ & 门长/ns\\\midrule\endfirsthead
\multicolumn{8}{c}{续表：R6 down：19层，769.396679 ns}\\\toprule
序 & 层 & 起点/ns & 门 & 线 & $\phi_{\rm nom}$ & $\phi_{\rm buf}$ & 门长/ns\\\midrule\endhead
\midrule\multicolumn{8}{r}{下页续}\endfoot
\bottomrule\endlastfoot
1 & 1 & 0.000 & $X_{01}/2$ & L & -1.57080 & -1.57080 & 32.000\\
2 & 2 & 32.000 & $\CZ$ & I,L & --- & --- & 56.371\\
3 & 3 & 88.371 & $X_{01}/2$ & L & 1.57080 & -2.58760 & 32.000\\
4 & 3 & 88.371 & $X_{01}$ & I & 0.00000 & 0.26136 & 32.000\\
5 & 4 & 120.371 & $X_{12}$ & I & 0.00000 & 0.24924 & 40.000\\
6 & 5 & 160.371 & $X_{01}/2$ & I & -1.57080 & -1.20202 & 32.000\\
7 & 6 & 192.371 & $\CZ$ & C,I & --- & --- & 30.642\\
8 & 7 & 223.013 & $X_{01}/2$ & I & 1.57080 & -1.23406 & 32.000\\
9 & 8 & 255.013 & $X_{12}$ & I & 0.00000 & -2.80294 & 40.000\\
10 & 9 & 295.013 & $X_{01}$ & I & 0.00000 & -2.93532 & 32.000\\
11 & 10 & 327.013 & $X_{01}/2$ & I & -1.57080 & 2.05326 & 32.000\\
12 & 11 & 359.013 & $\CZ$ & L,I & --- & --- & 56.371\\
13 & 12 & 415.384 & $X_{01}/2$ & I & 1.57080 & -0.78866 & 32.000\\
14 & 13 & 447.384 & $X_{12}$ & I & 0.00000 & -2.39392 & 40.000\\
15 & 13 & 447.384 & $X_{01}/2$ & L & -1.57080 & 2.67878 & 32.000\\
16 & 14 & 487.384 & $\CZ$ & I,L & --- & --- & 56.371\\
17 & 15 & 543.756 & $X_{01}/2$ & L & 1.57080 & 1.66197 & 32.000\\
18 & 15 & 543.756 & $X_{01}/2$ & R & -1.57080 & -1.57080 & 32.000\\
19 & 16 & 575.756 & $\CZ$ & I,R & --- & --- & 64.820\\
20 & 17 & 640.576 & $X_{01}/2$ & R & 1.57080 & 0.44708 & 32.000\\
21 & 17 & 640.576 & $X_{01}/2$ & I & -1.57080 & 2.74138 & 32.000\\
22 & 18 & 672.576 & $\CZ$ & R,I & --- & --- & 64.820\\
23 & 19 & 737.397 & $X_{01}/2$ & I & 1.57080 & -0.10435 & 32.000\\
24 & 19 & 737.397 & $F$ & C & 2.13903 & --- & 0.000\\
25 & 19 & 737.397 & $F$ & L & -0.13345 & --- & 0.000\\
26 & 19 & 737.397 & $F$ & R & 2.95696 & --- & 0.000\\
\end{longtable}}

\clearpage{\small\setlength{\tabcolsep}{4pt}\fontsize{9}{10.6}\selectfont\renewcommand{\arraystretch}{1.0}
\begin{longtable}{@{}rrrllrrr@{}}
\caption{R6 up：19层，769.396679 ns}\label{tab:r6up}\\\toprule
序 & 层 & 起点/ns & 门 & 线 & $\phi_{\rm nom}$ & $\phi_{\rm buf}$ & 门长/ns\\\midrule\endfirsthead
\multicolumn{8}{c}{续表：R6 up：19层，769.396679 ns}\\\toprule
序 & 层 & 起点/ns & 门 & 线 & $\phi_{\rm nom}$ & $\phi_{\rm buf}$ & 门长/ns\\\midrule\endhead
\midrule\multicolumn{8}{r}{下页续}\endfoot
\bottomrule\endlastfoot
1 & 1 & 0.000 & $F$ & C & 2.11427 & --- & 0.000\\
2 & 1 & 0.000 & $F$ & L & -0.14382 & --- & 0.000\\
3 & 1 & 0.000 & $F$ & R & 3.00675 & --- & 0.000\\
4 & 1 & 0.000 & $X_{01}/2$ & I & -1.57080 & -1.57080 & 32.000\\
5 & 2 & 32.000 & $\CZ$ & R,I & --- & --- & 64.820\\
6 & 3 & 96.820 & $X_{01}/2$ & I & 1.57080 & 1.86666 & 32.000\\
7 & 3 & 96.820 & $X_{01}/2$ & R & -1.57080 & 0.27602 & 32.000\\
8 & 4 & 128.820 & $\CZ$ & I,R & --- & --- & 64.820\\
9 & 5 & 193.641 & $X_{01}/2$ & R & 1.57080 & 2.29389 & 32.000\\
10 & 5 & 193.641 & $X_{01}/2$ & L & -1.57080 & -1.71461 & 32.000\\
11 & 6 & 225.641 & $\CZ$ & I,L & --- & --- & 56.371\\
12 & 7 & 282.012 & $X_{01}/2$ & L & 1.57080 & -2.73142 & 32.000\\
13 & 7 & 282.012 & $X_{12}$ & I & 0.00000 & 0.80093 & 40.000\\
14 & 8 & 322.012 & $X_{01}/2$ & I & -1.57080 & -0.88649 & 32.000\\
15 & 9 & 354.012 & $\CZ$ & L,I & --- & --- & 56.371\\
16 & 10 & 410.384 & $X_{01}/2$ & I & 1.57080 & 2.55477 & 32.000\\
17 & 11 & 442.384 & $X_{01}$ & I & 0.00000 & 1.02229 & 32.000\\
18 & 12 & 474.384 & $X_{12}$ & I & 0.00000 & 1.20994 & 40.000\\
19 & 13 & 514.384 & $X_{01}/2$ & I & -1.57080 & -0.44109 & 32.000\\
20 & 14 & 546.384 & $\CZ$ & C,I & --- & --- & 30.642\\
21 & 15 & 577.025 & $X_{01}/2$ & I & 1.57080 & -0.47313 & 32.000\\
22 & 16 & 609.025 & $X_{12}$ & I & 0.00000 & -1.84223 & 40.000\\
23 & 17 & 649.025 & $X_{01}$ & I & 0.00000 & -2.17439 & 32.000\\
24 & 17 & 649.025 & $X_{01}/2$ & L & -1.57080 & 2.53496 & 32.000\\
25 & 18 & 681.025 & $\CZ$ & I,L & --- & --- & 56.371\\
26 & 19 & 737.397 & $X_{01}/2$ & L & 1.57080 & 1.51815 & 32.000\\
\end{longtable}}


\section{工程组织与复现}\label{app:reproduce}
主文件\code{main.tex}，正文分为\code{sections/}内五个源文件；\code{images/}含25张矢量图，\code{tables/}含JSON导出的数表，\code{data/}保留主要结果与冻结参数。XeLaTeX两遍可解析目录与引用，无需图片公式。

\code{supplement/}保存R5/R6完整整网逐层Markdown及调度JSON，另有局域名义与缓冲相位JSON。主文宏观时间表和局域门表用于讲清逻辑，数百行完整执行表单独提供以便核对。

上一级\code{prepare_latex_report.py}只导出数表、复制图及记录指纹，不下发实验、不重跑传播。正文源文件不被该脚本覆盖。完整数值入口、依赖和原始报告见上一级\code{README.md}。本版保留旧PDF，单独提供原生公式和重构叙述。

\end{document}
